\documentclass[11pt,a4paper]{article}

\usepackage[utf8]{inputenc}
\usepackage[T1]{fontenc}
\usepackage{polski}
\usepackage{graphicx} 
\usepackage{geometry}
\geometry{margin=2.5cm}
\usepackage{setspace}
\onehalfspacing
\usepackage{amsmath, amssymb}
\usepackage[hidelinks]{hyperref}
\usepackage{amsthm}
\usepackage{algorithm}
\usepackage{algpseudocode}
\usepackage{booktabs}

\makeatletter
\renewcommand{\ALG@name}{Algorytm}
\makeatother
\theoremstyle{definition} 
\newtheorem{definition}{Definicja}[section]

\title{Totalne kolorowanie grafu - Raport}
\author{Klaudia Buczek, Klara Pluta, Gabriela Warchoł}
\date{}

\begin{document}

\maketitle

\section{Opis problemu}
Totalne kolorowanie grafu to problem polegający na przypisaniu kolorów zarówno wierzchołkom, jak i krawędziom grafu, tak aby żadne dwa sąsiednie elementy (wierzchołki lub krawędzie) nie miały tego samego koloru. Formalnie, totalne kolorowanie grafu możemy opisać przez funckję: $f: V \cup E \to \{1, 2, \ldots, k\}$, gdzie $V$ to zbiór wierzchołków, $E$ to zbiór krawędzi, a $k$ to liczba dostępnych kolorów. Ponadto muszą zachodzić następujące warunki:
\begin{itemize}
  \item Dla każdej krawędzi $uv \in E$ zachodzi $f(u) \neq f(v)$.
  \item Dla każdych dwóch różnych krawędzi $uv, ux \in E$ zachodzi $f(uv)\neq f(ux)$.
  \item Dla każdej krawędzi $uv \in E$ zachodzi $f(u) \neq f(uv)$ oraz $f(v) \neq f(uv)$.
\end{itemize}
Celem projektu jest wykorzystanie SAT-solverów do zbadania, czy dla zadanego grafu $G$ oraz parametru $k\in \mathbb{N}$ istnieje totalne kolorowanie.

\section{Modelowanie formuły logicznej}
Aby rozwiązać problem totalnego kolorowania grafu za pomocą SAT-solvera, musimy przekształcić go w formułę logiczną. W tym celu wprowadzamy zmienne logiczne, które reprezentują przypisanie kolorów do wierzchołków i krawędzi grafu.

\subsection{Zmienne logiczne}
Definiujemy zmienną $x_{a,b}$ dla każdego elementu $a \in V \cup E$ oraz każdego koloru $b \in K$ (gdzie $K$ to zbiór dostępnych kolorów). Zmienna $x_{a,b}$ przyjmuje wartość 1, jeśli element $a$ jest pokolorowany kolorem $b$, i 0 w przeciwnym przypadku.

\subsection{Ograniczenia}
Aby zapewnić poprawność totalnego kolorowania, musimy wprowadzić odpowiednie ograniczenia:
\begin{itemize}
  \item Każdy element (wierzchołek lub krawędź) musi być pokolorowany dokładnie jednym kolorem. To ograniczenie dzielimy na dwie części:
  \begin{itemize}
    \item Przynajmniej jeden kolor: $\forall a \in V \cup E \quad (x_{a,1} \lor x_{a,2} \lor \ldots \lor x_{a,k})$
    \item Co najwyżej jeden kolor: $\forall a \in V\cup E \quad \forall b, c \in K, b < c \quad (\neg x_{a,b} \lor \neg x_{a,c})$
  \end{itemize} 
  \item Sąsiednie wierzchołki nie mogą mieć tego samego koloru:
  \[
  \forall uv \in E, b \in K \quad (\neg x_{u,b} \lor \neg x_{v,b})
  \]
  \item Krawędzie wychodzące z jednego wierzchołka muszą mieć różne kolory:
  \[
  \forall e_1, e_2 \in E: e_1 \cap e_2 \neq \emptyset \quad \forall b \in K \quad (\neg x_{e_1, b} \lor \neg x_{e_2,b})
  \]
  \item Krawędź i jej końce nie mogą mieć tego samego koloru:
  \[
  \forall uv \in E, b \in K \quad (\neg x_{u,b} \lor \neg x_{uv,b}) \land (\neg x_{v,b} \lor \neg x_{uv,b})
  \]
\end{itemize}

\section{Implementacja}
Program został podzielony na kilka modułów, z których każdy odpowiada za inną część procesu rozwiązywania problemu totalnego kolorowania grafu.
\begin{itemize}
  \item \texttt{total\_coloring.py} - Główny program, wczytuje plik z grafem, zmienia krawędzie, wierzchołki i kolory na zmienne logiczne oraz tworzy formułę logiczną w formacie DIMACS.
  \item \texttt{run\_experiments.py} - Skrypt do automatycznego uruchamiania testów. Bierze po kolei grafy z folderu \texttt{examples} i sprawdza je dla coraz większej liczby kolorów, zapisując wyniki w pliku \texttt{results.txt}. Skrypt zatrzymuje się od razu, gdy solver zwróci pierwszy wynik \texttt{SAT}, dzięki temu wiemy, jaka jest minimalna liczba kolorów dla danego grafu.
  \item \texttt{verify\_unsat.py} - Kod do wyciągania dowodów niespełnialności za pomocą solvera \texttt{CaDiCaL}. 
  \item \texttt{check\_cnf.py} - Skrypt do sprawdzenia jednego pliku CNF przez solver \texttt{Glucose3}.
\end{itemize}  
\subsection{Algorytm działania pętli testowej}
Działanie skryptu uruchamiającego eksperymenty pokazuje Algorytm \ref{alg:pętla}.

\begin{algorithm}[H]
\caption{Automatyczne testowanie grafów}\label{alg:pętla}
\begin{algorithmic}[1]
\Require Folder z grafami \textbf{Examples}, zakres kolorów $K\_VALUES = [2, \dots, 13]$
\Ensure Plik z wynikami \texttt{results.csv}
\For{graf $g \in \text{Examples}$}
    \State Wczytaj wierzchołki $V$ i krawędzie $E$
    \For{$k \in K\_VALUES$}
        \State Wygeneruj formułę CNF dla parametrów $(G, k)$
        \State Uruchom SAT-solver
        \State Zapisz wyniki (liczbę zmiennych, klauzul i wynik SAT/UNSAT)
        \If{$result == \text{"SAT"}$}
            \State \textbf{break} \Comment{Mamy wynik, nie liczymy dla większych $k$}
        \EndIf
    \EndFor
\EndFor
\end{algorithmic}
\end{algorithm}

\section{Wyniki eksperymentów}
Przetestowałyśmy nasz model na wielu różnych rodzajach grafów (ścieżki, cykle, grafy pełne, dwudzielne, koła, gwiazdy oraz graf Petersena). 

\subsection{Minimalna liczba kolorów}
W tabeli \ref{tab:kolory} widać, przy jakiej najmniejszej liczbie kolorów $k$ solver był w stanie znaleźć poprawne totalne kolorowanie (czyli dał wynik \texttt{SAT}).

\begin{table}[H]
\centering
\caption{Wyniki minimalnego $k$ dla badanych grafów}\label{tab:kolory}
\begin{tabular}{lc|lc}
\toprule
Graf & Minimalne $k$ & Graf & Minimalne $k$ \\
\midrule
P4 & 3 & C10 & 4 \\
P6 & 3 & diamond & 4 \\
P8 & 3 & Q3 & 4 \\
P10 & 3 & K3,3 & 5 \\
C6 & 3 & K4 & 5 \\
C4 & 4 & K5 & 5 \\
C8 & 4 & K4,4 & 6 \\
W5 & 6 & Petersen & 6 \\
S5 & 6 & K5,5 & 7 \\
W7 & 8 & K6 & 7 \\
S8 & 9 & & \\
\bottomrule
\end{tabular}
\end{table}

Wszystkie wyniki zgadzają się z teorią grafów. Im graf ma więcej krawędzi i im więcej z nich schodzi się w jednym wierzchołku, tym więcej kolorów potrzebuje solver, żeby spełnić wszystkie warunki.

\subsection{Rozmiar wygenerowanych plików CNF}
Liczba zmiennych w formule zawsze zależy tylko od rozmiaru grafu i liczby kolorów, co idealnie potwierdza wzór:
\[
\text{Liczba zmiennych} = (|V| + |E|) \cdot k
\]
Dużo ciekawsza jest liczba klauzul, czyli ograniczeń. Jest ona uzależniona od stopnia wierzchołków. Tabela \ref{tab:rozmiary} pokazuje największe formuły uzyskane w testach.

\begin{table}[H]
\centering
\caption{Największe formuły uzyskane w testach}\label{tab:rozmiary}
\begin{tabular}{lcccc}
\toprule
Graf & $k$ & Zmienne & Klauzule & Wynik \\
\midrule
K5,5 & 7 & 245 & 1995 & SAT \\
W7 & 8 & 176 & 1310 & SAT \\
K6 & 7 & 147 & 1197 & SAT \\
S8 & 9 & 153 & 1097 & SAT \\
\bottomrule
\end{tabular}
\end{table}

Gdy graf ma dużo krawędzi, liczba klauzul drastycznie rośnie. Dla prostej ścieżki $P_{10}$ program wygenerował tylko 181 klauzul, a dla grafu pełnego $K_6$ o podobnej liczbie elementów (wierzchołków i krawędzi) wyszło ich aż 1197. 

\subsection{Weryfikacja przypadków UNSAT i DRAT}
W projekcie spróbowałyśmy uruchomić pełną weryfikację za pomocą programu \texttt{drat-trim}. Udało nam się go skompilować i przygotować całe środowisko. Pojawił się jednak problem techniczny: biblioteka \texttt{PySAT}, której używamy w Pythonie do obsługi solverów, nie potrafiła wygenerować poprawnych, niepustych plików z dowodami (\textit{proofs}).

Z tego powodu nie mogłyśmy w pełni domknąć automatycznej weryfikacji przez \texttt{drat-trim}. Wszystkie przypadki \texttt{UNSAT} (gdy kolorów było za mało) zweryfikowałyśmy jednak bezpośrednio w solverze \texttt{Glucose3}. Każdy test dla mniejszego $k$ kończył się twardym wynikiem \texttt{UNSAT}, co udowadnia, że generator działa poprawnie i nie przepuszcza błędnych kolorowań.

\end{document}